\documentclass[10pt,a4paper]{amsart}
\usepackage[margin=19mm]{geometry}
\usepackage{amsmath,amssymb,mathtools}
\usepackage[scr=rsfs,cal=euler]{mathalfa}
\usepackage{xfrac}
\usepackage[unicode,hidelinks,bookmarks=false]{hyperref}
\newcommand{\ie}{i.e.\ }
\newcommand{\cf}{cf.\ }
\newcommand{\resp}{respectively\ }
\newcommand{\Subsection}{\subsection}
\providecommand{\nobreakdash}{\nobreak\hbox{-}}
\setlength{\emergencystretch}{2em}
\multlinegap0pt
\usepackage{bm}
\usepackage{amssymb}
\usepackage{xcolor}
\newtheorem{thm}{Theorem}
    \newcommand{\Hom}{{\mathrm{Hom}}} \renewcommand{\Im}{{\mathrm{Im}}}
\title{A note on small theta lift}
\author{Jingsong Chai}
\date{Source: arXiv:2604.10421v1 (CC BY 4.0)}
\begin{document}
\maketitle

\noindent\emph{Source and licence.} A note on small theta lift. Version arXiv:2604.10421v1 (CC BY 4.0), distributed by arXiv under the Creative Commons Attribution 4.0 International licence, http://creativecommons.org/licenses/by/4.0/, as recorded in the arXiv licence metadata for this exact version and shown on the abstract page https://arxiv.org/abs/2604.10421v1. Full licence notice: \texttt{licenses/arxiv-2604.10421-CC-BY-4.0.txt}.\par\smallskip

\noindent\textit{Editorial scope.} Counted unit: the article's thm thm (source0.tex lines 450--455). The article's complete original proof of it is delivered with it (source0.tex lines 456--504, one \texttt{proof} environment of 49 lines). No mathematics is written, repaired or re-derived: the statement and the proof are the article's own lines, in the article's own notation, and no retained line is substituted or re-flowed.  No further source-local context is needed: every cross-reference of every retained line resolves inside the two delivered spans.  Nothing was omitted from any retained span, so the package carries no omission record.  The retained lines cite 1 works, whose entries are transcribed verbatim from the article's own bibliography source in the e-print and delivered with short editorial labels recorded in the item's reference mapping.  No retained line contains a figure environment, an image-inclusion command or drawing code, so no image is delivered and the package declares no asset dependency.  Editorial formatting only, in addition: an AMS \texttt{amsart} preamble that restates the article's own declarations and macros used by the retained lines, namely the environments \texttt{thm} on separate counters, together with the standard packages those lines need.  The retained environments print in this document as Theorem 1 in order of appearance, whereas the article prints the same items with the numbering of its own full document.  The complete native source of the article as distributed in the arXiv e-print for this exact version is bundled unchanged as \texttt{sources/198273/source0.tex}.
\begin{thm}
We have 
\[
H_{\ell,\psi}(\pi)\cong \theta_{V,W,\boldsymbol{\chi},\psi}(\pi).
\]
\end{thm}

\begin{proof}
Since 
\[
\Theta_{V,W,\boldsymbol{\chi},\psi}(\pi)=(\omega_{V,W,\boldsymbol{\chi},\psi}\otimes \pi^\vee)_{G(V)},
\]
it follows that $H_{\ell,\psi}(\pi)$ is a quotient of $\Theta_{V,W,\boldsymbol{\chi},\psi}(\pi)$. Hence it is of finite length. By the Howe duality, $\Theta_{V,W,\boldsymbol{\chi},\psi}(\pi)$ has a unique quotient $\theta_{V,W,\boldsymbol{\chi},\psi}(\pi)$. Thus $\theta_{V,W,\boldsymbol{\chi},\psi}(\pi)$ is also the unique quotient of $H_{\ell,\psi}(\pi)$. Hence we can write
\[
\theta_{V,W,\boldsymbol{\chi},\psi}(\pi) = (\omega_{V,W,\boldsymbol{\chi},\psi}\otimes \pi^\vee)/R_1,
\]
where $R\subseteq R_1\subset \omega_{V,W,\boldsymbol{\chi},\psi}\otimes \pi^\vee$.

If $H_{\ell,\psi}(\pi)$ is not irreducible, then $R_1$ will contain $R$ properly. Now consider
\[
H_{\ell, \bar{\psi}}(\pi^\vee \chi_V):=(\bar{\omega}_{V,W,\boldsymbol{\chi},\psi} \otimes \pi \chi_V^{-1} )/R',
\]
where $R'\subset \bar{\omega}_{V,W,\boldsymbol{\chi},\psi} \otimes \pi \chi_V^{-1}$ is defined similarly as $R$. Also $\theta_{V,W,\boldsymbol{\chi},\bar{\psi}}(\pi^\vee \chi_V)$ is a quotient of $H_{\ell,\bar{\psi}}(\pi^\vee \chi_V)$ and can be written as 
\[
\theta_{V,W,\boldsymbol{\chi},\bar{\psi}}(\pi^\vee \chi_V)=(\bar{\omega}_{V,W,\boldsymbol{\chi},\psi} \otimes \pi \chi_V^{-1} )/R_1'
\]
with $R'\subseteq R_1' \subset \bar{\omega}_{V,W,\boldsymbol{\chi},\psi} \otimes \pi \chi_V^{-1} $.

By Corollary 18.4 in \cite{GKT}, there is a natural nonzero linear form in 
\[
\Hom_{G(V)}(\theta_{V,W,\boldsymbol{\chi},\psi}(\pi)\otimes \theta_{V,W,\boldsymbol{\chi},\bar{\psi}}(\pi^\vee \chi_V), \chi_W).
\]
View this nonzero linear form as an element in 
\[
\Hom_{G(V)}((\omega_{V,W,\boldsymbol{\chi},\psi}\otimes \pi^\vee)/R_1 \otimes (\bar{\omega}_{V,W,\boldsymbol{\chi},\psi} \otimes \pi \chi_V^{-1} )/R_1' ,\chi_W),
\]
which can be naturally extended to a nonzero liner form in 
\[
\Hom_{G(V)}( H_{\ell,\psi}(\pi)\otimes H_{\ell,\bar{\psi}}(\pi^\vee \chi_V),\chi_W ).
\]
Extend this nonzero linear form further to an element in
\[
\Hom_{G(V)}( \Theta_{V,W,\boldsymbol{\chi},\psi}(\pi)\otimes \Theta_{V,W,\boldsymbol{\chi},\bar{\psi}}(\pi^\vee \chi_V), \chi_W  ).
\]
This nonzero element will correspond to a nonzero element $\ell'$ in
\[
\Hom_{G'(W)\times G'(W)\times G(V)} (\omega_{V,W,\boldsymbol{\chi},\psi}\otimes \bar{\omega}_{V,W,\boldsymbol{\chi},\psi} \otimes \pi \chi_V^{-1}\otimes \pi^\vee , \chi_W ).
\]
As this Hom space has dimention at most one, there exists a nonzero scalar $\lambda$, such that
\[
\ell=\lambda \ell'.
\]
But, by the construction of $\ell'$, the radical is $R_1$, which is strictly larger than $R$, and we got a contradiction, which will imply that $H_{\ell,\psi}(\pi)$ is irreducible.

So $H_{\ell,\psi}(\pi)$ is an irreducible quotient of $\Theta_{V,W,\boldsymbol{\chi},\psi}(\pi)$, it has to be $\theta_{V,W,\boldsymbol{\chi},\psi}(\pi)$ by Howe duality, and the conclusion follows.
\end{proof}

\begin{thebibliography}{99}
\bibitem[GKT]{GKT} Wee Tech Gan, S. Kudla and S. Takeda. {\em The local theta correspondence.} Preprint.
\end{thebibliography}
\end{document}
